% problem-000186 8 烯烃光氧化裂解
\section*{8. 烯烃光氧化裂解}
实验⼀：化合物1在 $\operatorname { C u C l } ( \operatorname { c o d } )$ _{2}催化下发⽣反应，⾼区域选择性形成产物2，其中2a和2b的⽐例为98:2，产物2a中的双键以反式为主，反式和顺式的⽐例为74:26。

（图：）

实验⼆：该反应在没有Cu(I)参与下，基本上不反应，原料完全回收。

实验三：产物(�)-2a可以在没有Cu(I)参与下，在100 °C下转化为(�)-2a和(�)-2b。然⽽，在同样条件下，

(�)-2a 则保持不变，不会转化为(�)-2a。

（图：）

分析以上信息，回答相关问题:

\subsection*{8-1}
为实验⼀转化为2a的过程提供关键中间体；

\subsection*{8-2}
为实验三(�)-2a异构化为(�)-2a、形成(�)-2b的过程提供关键中间体；

\subsection*{8-3}
在加热条件下，解释为什么(�)-2a 可以转化为(�)-2a，⽽(�)-2a 不能转化为(�)-2a；

\subsection*{8-4}
在后续的研究过程中，发现除了形成以上四元环外，还可以形成五元杂环 $3 _ { \circ }$ 。画出形成化合物3的关键中间体（说明：与前⾯重复的⽆需再画）。

（图：）
